% ============================================================
%  Оформление отчёта. Обычно здесь ничего менять не нужно.
%  Сборка по умолчанию выполняется XeLaTeX.
% ============================================================

\usepackage{iftex}
\ifPDFTeX
  \usepackage[T2A]{fontenc}
  \usepackage[utf8]{inputenc}
\else
  \usepackage{fontspec}
  \setmainfont{Liberation Serif}
  \setmonofont{Liberation Mono}[Scale=MatchLowercase]
\fi
\usepackage[english,russian]{babel}

% --- Поля и абзацы (как в исходном doc-шаблоне курса) ---
\usepackage[a4paper, margin=2.54cm]{geometry}
\setlength{\parindent}{0pt}
\setlength{\parskip}{4pt plus 1pt}
\frenchspacing

% --- Заголовки разделов без номеров ---
\usepackage{titlesec}
\setcounter{secnumdepth}{0}
\titleformat{\section}{\large\bfseries}{}{0pt}{}
\titleformat{\subsection}{\normalsize\bfseries}{}{0pt}{}
\titlespacing*{\section}{0pt}{14pt plus 4pt minus 2pt}{4pt plus 2pt}
\titlespacing*{\subsection}{0pt}{10pt plus 2pt minus 2pt}{3pt plus 1pt}

% --- Таблицы, рисунки, подписи ---
\usepackage[table]{xcolor}
\usepackage{graphicx}
\usepackage{tabularx}
\usepackage{float}
\usepackage{caption}
\DeclareCaptionLabelSeparator{spacedendash}{~-- }
\captionsetup{labelsep=spacedendash, font=small, singlelinecheck=false, justification=raggedright}
\captionsetup[table]{position=top, skip=4pt}
\DeclareCaptionType[fileext=lol]{listing}[Листинг][Листинги]
\captionsetup[listing]{position=top, skip=4pt, hypcap=false}
\addto\captionsrussian{%
  \renewcommand{\figurename}{Рисунок}%
  \renewcommand{\tablename}{Таблица}%
}
\newcommand{\thead}[1]{\multicolumn{1}{|c|}{\cellcolor{gray!20}\textbf{#1}}}
% Колонки таблиц без растягивания пробелов: L – резиновая, P{ширина} – фиксированная
\newcolumntype{L}{>{\raggedright\arraybackslash}X}
\newcolumntype{P}[1]{>{\raggedright\arraybackslash}p{#1}}

% --- Код и вывод терминала ---
% code     – фрагмент исходного кода;
% terminal – вывод программы, strace и т. п.;
% \inputterminal{файл} / \inputcode{файл} – то же, но из файла.
\usepackage{fvextra}
\fvset{breaklines=true, breakanywhere=true,
       frame=leftline, framerule=1.5pt, framesep=6pt,
       rulecolor=\color{gray!60}, xleftmargin=4pt}
\DefineVerbatimEnvironment{code}{Verbatim}{fontsize=\small}
\DefineVerbatimEnvironment{terminal}{Verbatim}{fontsize=\footnotesize}
\newcommand{\inputcode}[2][]{\VerbatimInput[fontsize=\small,#1]{#2}}
\newcommand{\inputterminal}[2][]{\VerbatimInput[fontsize=\footnotesize,#1]{#2}}
% Листинг с подписью сверху:
%   \begin{listingblock}{Подпись} \begin{code}...\end{code} \end{listingblock}
\newenvironment{listingblock}[1]
  {\par\medskip\captionof{listing}{#1}\nopagebreak\vspace{-\topsep}\vspace{-\parskip}}
  {\par}

% Приложение: с новой страницы, заголовок «Приложение А. ...»,
% листинги внутри нумеруются А.1, А.2, ...
\newcounter{appendixno}
\newcommand{\appendixsection}[2]{%
  \clearpage
  \section{Приложение #1. #2}%
  \stepcounter{appendixno}%
  \setcounter{listing}{0}%
  \renewcommand{\thelisting}{#1.\arabic{listing}}%
  \renewcommand{\theHlisting}{app\arabic{appendixno}.\arabic{listing}}}

% --- Ссылки ---
\usepackage[hidelinks]{hyperref}
\usepackage{xurl}

% --- Подсказки шаблона: серый курсив, удалить перед сдачей ---
\newcommand{\hint}[1]{\leavevmode{\color{gray}\itshape #1}}
% Рамка-заглушка вместо рисунка
\newcommand{\placeholder}[2][4.5cm]{%
  \fbox{\parbox[c][#1][c]{0.95\linewidth}{\centering\hint{#2}}}}
